\documentclass[11pt,a4paper]{article}
\usepackage[utf8]{inputenc}
\usepackage[T1]{fontenc}
\usepackage[spanish,es-nodecimaldot]{babel}
\usepackage{lmodern}
\usepackage[top=2.3cm,bottom=2.3cm,left=2.4cm,right=2.4cm]{geometry}
\usepackage{booktabs,tabularx,array,enumitem,amsmath,hyperref}
\usepackage{pgfplots}
\pgfplotsset{compat=1.18}
\hypersetup{colorlinks=true,urlcolor=blue,linkcolor=black}
\setlength{\parindent}{0pt}
\setlength{\parskip}{0.5em}
\renewcommand{\arraystretch}{1.14}
\setlist[itemize]{topsep=3pt,itemsep=2pt,parsep=0pt}
\urlstyle{tt}
\newcolumntype{Y}{>{\raggedright\arraybackslash}X}
\newcommand{\archivo}[1]{\nolinkurl{#1}}

\title{\textbf{Exploración de modelos de Inteligencia Artificial multimodal\\
para la identificación de macromicetos mediante dispositivos móviles}\\[0.6em]
\normalsize Semanas 3 y 4: curaduría del dataset y formulación de la línea base visual con ResNet-18}
\author{Daniel Esteban Camacho Ospina\\Cesar Augusto Pulido Cuervo}
\date{4 de octubre de 2026}

\begin{document}
\maketitle

\begin{abstract}
Se establece una línea base visual para la clasificación de diez géneros de
macromicetos mediante transferencia de aprendizaje con ResNet-18. El dataset
v002 comprende 43\,621 imágenes de 25\,290 observaciones de FungiTastic y conserva
las particiones oficiales de entrenamiento, validación y prueba. La metodología
incluye normalización taxonómica, un esquema de procedencia y calidad de los
metadatos ambientales, versionado y controles de integridad y fuga de información.
La selección del modelo utiliza macro-F1 por observación en validación.
El entrenamiento concluye en la época 16 mediante parada temprana y selecciona
el checkpoint de la época 11, con macro-F1 de observación de 0,9545 en val y
0,9463 en test. La evaluación final de test se ejecutó una sola vez.
El trabajo corresponde a las semanas S3 y S4 y
establece las condiciones para evaluar una arquitectura eficiente para móvil
y consolidar las métricas durante S5 y S6.
\end{abstract}

\section{Introducción y objetivo}

La identificación de macromicetos a partir de fotografías requiere una base de
datos consistente y un protocolo de evaluación que permita distinguir el
aprendizaje visual de la dependencia entre observaciones. La coexistencia de
varias imágenes de un mismo hallazgo, el desbalance entre géneros y las diferencias
de contexto geográfico condicionan la interpretación del desempeño.

El objetivo de esta etapa es establecer un dataset reproducible y formular una
arquitectura de referencia para cuantificar el rendimiento inicial del componente
visual. El modelo \(M_A=f(I)\) utiliza únicamente la imagen \(I\) para predecir uno
de diez \textbf{géneros}. Su alcance es la clasificación a nivel de género;
la identificación de especies y la determinación de comestibilidad no son
variables objetivo de este experimento. Aunque el proyecto contempla modelos
multimodales, los metadatos ambientales se conservan para etapas posteriores y
no se incorporan a la entrada de esta línea base.

\section{Correspondencia con las semanas S3 y S4}

\begin{center}
\small
\begin{tabularx}{\textwidth}{p{2.7cm}p{3.1cm}Y}
\toprule
\textbf{Semana y periodo} & \textbf{Fase} & \textbf{Alcance de la etapa}\\
\midrule
S3: 21--27 de septiembre & Fase 2: Auditoría y curaduría & Reglas de limpieza y
normalización; particiones reproducibles train/val/test; esquema de metadatos
ambientales con fuente, método y calidad; versionado del dataset.\\
S4: 28 de septiembre--4 de octubre & Fase 3: Línea base visual & Configuración
del entorno de entrenamiento y primer entrenamiento con una arquitectura de referencia.\\
\bottomrule
\end{tabularx}
\end{center}

El dataset v002 y su protocolo de particiones materializan los componentes de
curaduría y reproducibilidad de S3. La configuración de ResNet-18 y el
entrenamiento inicial corresponden a S4. La arquitectura de referencia completó
su evaluación final y sus artefactos quedaron conservados. El análisis distingue
el criterio de selección en val de las métricas finales de test.

\section{Diseño del dataset y trazabilidad}

La versión v002 utiliza una sola fuente y preserva las particiones oficiales de
FungiTastic. El versionado distingue esta definición de versiones anteriores y
permite asociar cada experimento con un conjunto específico de imágenes,
etiquetas y controles de calidad. La integridad se comprueba con SHA-256 y la
consistencia se valida entre las tablas de observaciones e imágenes.

\subsection{Organización y trazabilidad}

\begin{center}
\small
\begin{tabularx}{\textwidth}{p{5.7cm}Y}
\toprule
\textbf{Ubicación} & \textbf{Contenido y función}\\
\midrule
\archivo{data/raw/fungitastic/} & Metadatos originales; se leen sin modificarlos.\\
\archivo{data/raw/images/fungitastic_v002/} & Fotografías locales y manifiesto de descarga.\\
\archivo{data/interim/v002/} & Tablas candidatas necesarias para planificar la descarga.\\
\archivo{data/processed/v002/} & Tablas finales \archivo{observations.csv} e \archivo{images.csv}.\\
\archivo{data/versions/v002/} & Manifiesto, clases, particiones, auditorías y ficha del dataset.\\
\archivo{configs/train/ma_resnet18_v002.yaml} & Arquitectura, transformaciones y parámetros de entrenamiento.\\
\archivo{runs/<run_id>/} & Configuración efectiva, historial, checkpoints, progreso y resultados.\\
\bottomrule
\end{tabularx}
\end{center}

Las fotografías, los datos procesados y las salidas de entrenamiento permanecen en
local y se excluyen de Git. Las configuraciones y los artefactos de auditoría se
pueden versionar. El manifiesto registra hashes de entradas, código y salidas.
La reconstrucción de v002 con \archivo{--check} produjo contenido idéntico.
La huella publicada es:
\begin{center}\small\ttfamily
144323e165ecfbe86d29a3484e56fd91b8e8cb1dce1eef5f1d1355fea7bedb90
\end{center}

\subsection{Esquema de metadatos ambientales}

La tabla de observaciones organiza la información ambiental en grupos de
geolocalización, fecha, elevación, sustrato, hábitat, cobertura del suelo, región
biogeográfica y clima. Cada grupo asocia sus variables con una terna de
procedencia: \textbf{fuente}, \textbf{método} y \textbf{calidad}.

La fuente identifica el origen del dato; el método distingue la forma de
obtención o transformación; la calidad permite separar valores válidos,
imprecisos, sospechosos, ausentes o invalidados. Esta representación evita
interpretar una ausencia como un valor observado y permite definir criterios
explícitos para futuros modelos multimodales. La normalización conserva además
el texto original cuando se utilizan vocabularios controlados.

\section{Fuente, curaduría y tamaño final}

\subsection{Selección de FungiTastic}

Se emplean exclusivamente los CSV \archivo{Train}, \archivo{ClosedSet-Val}
y \archivo{ClosedSet-Test}. No se creó \archivo{test_external}.
La auditoría previa y la configuración documentan la exclusión de DF20 por su
solapamiento con FungiTastic, de OpenSet porque contiene ClosedSet, y de
Mini/FewShot y DNA por sus protocolos distintos. GBIF, iNaturalist,
Mushroom Observer y MIND.Funga no forman parte de esta primera versión de una fuente.
Estas decisiones reducen el riesgo de duplicar contenido entre fuentes; no demuestran
por sí solas ausencia de sesgo.

\subsection{Unidad de análisis y limpieza}

Una \textbf{observación} representa un hallazgo y puede contener varias fotografías.
Cada imagen tiene un identificador propio y una referencia a su observación.
Todas las imágenes del mismo hallazgo conservan la etiqueta y la partición de este.

La lectura de los tres CSV produjo 615\,193 filas de imagen. El registro de limpieza
documenta la exclusión de 10\,049 filas sin clasificación en Fungi y, posteriormente,
561\,523 filas fuera de los diez géneros. Quedaron \textbf{43\,621 imágenes de
25\,290 observaciones}. Estos conteos corresponden al subconjunto seleccionado,
no a la totalidad de FungiTastic.

Se normalizan nombres, espacios y sinónimos configurados
(\textit{Coriolus}\(\rightarrow\)\textit{Trametes},
\textit{Hirneola}\(\rightarrow\)\textit{Auricularia}); los identificadores de
clase van de 0 a 9 según el orden de la configuración. También se normalizan
fechas, coordenadas, licencias y vocabularios ambientales, conservando su procedencia
y calidad. Las marcas de calidad ambiental no implican excluir una imagen de este
modelo. Por ejemplo, el clima figura como ausente en las 25\,290 observaciones.

\subsection{Conteos por género}

Los conteos de imágenes y observaciones son números enteros; los millares se
separan mediante un espacio fino. Los porcentajes se identifican explícitamente
con el símbolo \(\%\), y la coma se utiliza como separador decimal.

\begin{center}
\small
\begin{tabular}{lrrr|rrr}
\toprule
& \multicolumn{3}{c}{\textbf{Observaciones}} & \multicolumn{3}{c}{\textbf{Imágenes}}\\
\textbf{Género} & \textbf{Train} & \textbf{Val} & \textbf{Test} & \textbf{Train} & \textbf{Val} & \textbf{Test}\\
\midrule
Amanita & 3\,752 & 904 & 1\,315 & 6\,508 & 1\,709 & 2\,393\\
Auricularia & 1\,443 & 288 & 319 & 2\,106 & 411 & 475\\
Ganoderma & 1\,832 & 335 & 346 & 3\,567 & 738 & 653\\
Geastrum & 779 & 79 & 182 & 1\,350 & 157 & 377\\
Hericium & 494 & 66 & 51 & 826 & 114 & 101\\
Lycoperdon & 2\,166 & 408 & 407 & 3\,336 & 686 & 660\\
Morchella & 256 & 11 & 20 & 405 & 21 & 34\\
Pleurotus & 1\,465 & 251 & 319 & 2\,690 & 495 & 616\\
Trametes & 3\,654 & 756 & 791 & 6\,626 & 1\,421 & 1\,384\\
Xylaria & 1\,958 & 354 & 289 & 2\,758 & 538 & 466\\
\midrule
\textbf{Total} & \textbf{17\,799} & \textbf{3\,452} & \textbf{4\,039} & \textbf{30\,172} & \textbf{6\,290} & \textbf{7\,159}\\
\bottomrule
\end{tabular}
\end{center}

Se conserva el desbalance original del subconjunto oficial. No se igualaron los
conteos por clase, no se sobremuestrearon clases y no se aplicaron pesos de clase
en la función de pérdida. La macro-F1 permite dar igual importancia a los diez
géneros durante la selección, pero no elimina el desbalance de las actualizaciones
de entrenamiento. Morchella tiene solo 11 observaciones de val y 20 de test;
su rendimiento debe interpretarse con este soporte reducido.

\section{Particiones y controles de fuga de información}

\subsection{Protocolo oficial}

\begin{center}
\small
\begin{tabularx}{\textwidth}{lrrY}
\toprule
\textbf{Partición} & \textbf{Observaciones} & \textbf{Imágenes} & \textbf{Uso}\\
\midrule
Train & 17\,799 (70,38\,\%) & 30\,172 (69,17\,\%) & Actualizar los pesos del modelo.\\
ClosedSet-Val & 3\,452 (13,65\,\%) & 6\,290 (14,42\,\%) & Seleccionar el checkpoint y controlar la parada temprana.\\
ClosedSet-Test & 4\,039 (15,97\,\%) & 7\,159 (16,41\,\%) & Evaluación final del mejor checkpoint.\\
\bottomrule
\end{tabularx}
\end{center}

Cada celda de observaciones e imágenes presenta el conteo y, entre paréntesis,
su porcentaje respecto al total de 25\,290 observaciones o 43\,621 imágenes,
respectivamente.

La columna \archivo{split} se asigna directamente a partir del CSV oficial y
coincide con \archivo{split_temporal}. Las proporciones anteriores se calcularon
después del filtrado de géneros: no son porcentajes fijados mediante una nueva
partición aleatoria de \(70\,\%/15\,\%/15\,\%\). Val es la evaluación durante el desarrollo.
Test queda reservado y no interviene en gradientes, selección o parada temprana.

\subsection{Auditorías ejecutadas}

\begin{itemize}
\item \textbf{Identidad:} cero identificadores de imagen u observación duplicados
y cero observaciones con imágenes en más de una partición. El esquema valida
la coincidencia de etiqueta, partición y número de fotografías de cada observación.
\item \textbf{Sitio-día:} se agrupan observaciones por fecha y coordenadas redondeadas
a dos decimales. Si faltan datos válidos, se usa un grupo individual. Se encontraron
\textbf{cero grupos compartidos entre particiones}. Este es un diagnóstico:
no se reasignan filas y la celda angular no constituye una distancia física constante.
\item \textbf{Duplicados exactos:} se verificaron los SHA-256 de las imágenes y no
se encontraron copias exactas entre particiones; las exclusiones fueron cero.
La política implementada protege test sobre val y val sobre train, retirando solo
la copia de menor prioridad. Las copias dentro de la misma partición se conservan.
\item \textbf{Similitud perceptual:} se utiliza pHash de 64 bits y distancia de
Hamming menor o igual a 4. Se reportó \textbf{un par entre train y test}, de
Lycoperdon: \archivo{0-3336959327.JPG} y \archivo{2-4465862839.JPG}, con distancia 4.
Se conserva para respetar el protocolo; similitud perceptual no prueba identidad.
\item \textbf{Tamaño:} 33 fotografías decodificables tienen un lado menor de 128
píxeles. Se conservaron y reportaron; las transformaciones las redimensionan.
\end{itemize}

Los controles se registran en \archivo{image_integrity.csv},
\archivo{image_exclusions.csv}, \archivo{near_duplicates.csv},
\archivo{site_day_overlaps.csv} y \archivo{small_images.csv}.
El entrenamiento vuelve a comprobar los archivos contra los hashes publicados.

\section{Formulación de la línea base visual con ResNet-18}

\subsection{Fundamento y papel de la arquitectura de referencia}

ResNet-18 es una red convolucional de 18 capas con conexiones residuales.
En un bloque de este tipo, la salida combina una transformación aprendida con
la entrada mediante una conexión de identidad o una proyección compatible:
\[
h_{l+1}=\sigma\bigl(F(h_l;\theta_l)+S(h_l)\bigr),
\]
donde \(F\) representa la transformación del bloque, \(S\) la conexión residual
y \(\sigma\) la activación. Esta estructura facilita la propagación de la
información y de los gradientes durante la optimización.

El preentrenamiento en ImageNet aporta una inicialización de filtros visuales
que se adaptan al dominio de macromicetos mediante ajuste de toda la red. La
arquitectura se utiliza como referencia para estudiar la capacidad discriminativa
de la imagen y comparar posteriormente modelos aptos para dispositivos móviles.
Su elección no acredita por sí sola eficiencia de inferencia en un teléfono:
esa propiedad requiere medir tamaño, latencia y consumo en la etapa siguiente.

\subsection{Arquitectura y entrada}

Se utiliza \textbf{ResNet-18} con pesos \archivo{IMAGENET1K_V1} de torchvision.
La capa final se reemplaza por una capa lineal de diez salidas; se ajustan todos
los parámetros de la red. El modelo produce logits y, en evaluación, se aplica
softmax para obtener probabilidades por género.

Las imágenes se convierten a RGB y la entrada de la red es de
\(224\times224\) píxeles. En train se aplican recorte aleatorio redimensionado
con escala de área entre 0,5 y 1,0, volteo horizontal aleatorio y variación de color
(brillo, contraste y saturación: 0,2; tono: 0,02).
En val/test se redimensiona el lado menor a 224 píxeles y se aplica recorte central
de \(224\times224\), sin aumento aleatorio. En ambos casos se normaliza cada
canal con las medias \((0,485;0,456;0,406)\) y desviaciones
\((0,229;0,224;0,225)\) de ImageNet.

\subsection{Pérdida y optimización}

La función de pérdida es entropía cruzada con suavizado de etiquetas
\(\varepsilon=0,1\). Para la clase objetivo \(y\) y \(K=10\):
\[
q_c=(1-\varepsilon)\mathbf{1}[c=y]+\frac{\varepsilon}{K},
\qquad \mathcal{L}=-\sum_{c=1}^{K}q_c\log p_c.
\]
No se incorporan ponderaciones de clase. El optimizador es AdamW; la tasa de
aprendizaje se incrementa durante la primera época y después sigue un descenso
coseno, actualizado por lote.

\begin{center}
\small
\begin{tabularx}{\textwidth}{p{5.4cm}Y}
\toprule
\textbf{Parámetro} & \textbf{Valor implementado}\\
\midrule
Máximo de épocas / semilla & 30 / 42\\
Lote train / lote evaluación & 64 / 128\\
Optimizador / tasa inicial & AdamW / \(3\times10^{-4}\)\\
Decaimiento de pesos & 0,05\\
Scheduler / calentamiento & Coseno / una época\\
Suavizado de etiquetas & 0,1\\
Selección / parada temprana & Macro-F1 por observación en val / cinco épocas sin mejora\\
Carga de datos & Cuatro workers; recreados por época para la reanudación\\
Dispositivo de la corrida revisada & CPU; AMP desactivado en la configuración efectiva\\
Determinismo estricto & Desactivado; se fija la semilla y se guardan estados aleatorios\\
\bottomrule
\end{tabularx}
\end{center}

Los lotes de train se barajan. Con 30\,172 imágenes y lote 64 se realizan
471 lotes completos por época; el último lote incompleto de 28 imágenes se
descarta en esa época. Val/test no descartan el último lote. La selección no
usa la pérdida de train: el mejor checkpoint maximiza macro-F1 de observación en val.

\subsection{Métricas por imagen y por observación}

Además de macro-F1, se calculan exactitud, exactitud balanceada, top-1 y top-3,
precisión, sensibilidad y F1 por género, soporte y matriz de confusión.
Las predicciones de imágenes de un mismo hallazgo se agregan mediante promedio
de probabilidades:
\[
\bar p_{o,c}=\frac{1}{n_o}\sum_{i\in I_o}p_{i,c},
\qquad \hat y_o=\arg\max_c\bar p_{o,c}.
\]
Así, cada observación cuenta una vez en la métrica utilizada para seleccionar el
modelo, aunque tenga varias fotografías. Para cada género se calcula
\(F1_c=2P_cR_c/(P_c+R_c)\), y
\(\mathrm{macro\text{-}F1}=\frac{1}{10}\sum_{c=1}^{10}F1_c\).
Las clases sin predicción positiva reciben cero en las divisiones indefinidas.

Al terminar el entrenamiento se carga \archivo{best.pt} y se producen métricas
y predicciones de val y test. Las rutas \archivo{--smoke} y
\archivo{--eval-only} evalúan solo val; la primera comprueba la ejecución con
tres lotes y no constituye una estimación de rendimiento.

\section{Reproducibilidad y continuidad experimental}

Cada ejecución conserva la configuración efectiva, la huella del dataset,
el historial de optimización y los estados necesarios para continuar el
experimento. La verificación de disponibilidad e integridad precede al
entrenamiento, mientras que el seguimiento del progreso permite distinguir
la preparación de datos, las actualizaciones del modelo y la validación.

\archivo{best.pt} conserva el mejor modelo según val.
\archivo{last.pt} se guarda de forma atómica al completar una época y su validación:
incluye pesos actuales, optimizador, scheduler, GradScaler, estados aleatorios,
historial, mejor checkpoint y contador de parada temprana.
\archivo{--resume} continúa desde la siguiente época y crea un nuevo directorio
con el historial heredado y el campo \archivo{resumed_from}.

La opción \archivo{-PauseAfterEpoch} o un archivo \archivo{pause.request} permite
pausar al terminar la época sin evaluar test. Una interrupción a mitad de época
pierde ese avance parcial, pero conserva el último \archivo{last.pt} completo.
Antes de finalizar la primera época todavía no existe este checkpoint.
La reanudación exige la misma configuración, etiquetas, huella y tipo de dispositivo;
la semilla no garantiza resultados numéricamente idénticos entre hardware distinto.

Se rechaza reanudar una corrida terminada o que ya inició la evaluación de test.
\archivo{test_started.json} registra el comienzo de esa evaluación y
\archivo{completed.json} acredita una finalización exitosa.
Estos controles evitan una repetición automática de test en esa corrida.

\section{Resultados y discusión}

\subsection{Finalización y selección de la arquitectura de referencia}

La ejecución consolidada es \archivo{20261004T133123181686Z_ma_resnet18_v002},
que heredó el historial y los estados de entrenamiento de las reanudaciones previas.
Se completaron \textbf{16 épocas}; la parada temprana ocurrió después de cinco
épocas consecutivas (12--16) sin superar la macro-F1 de observación en val
de la \textbf{época 11}. Este checkpoint fue recuperado para la evaluación final.
El registro de finalización confirma \textbf{una evaluación de test} y la
existencia de historial, predicciones y métricas finales.

La configuración efectiva corresponde a Python 3.12.10, PyTorch 2.14.1+cpu
y torchvision 0.29.1+cpu, con ejecución en CPU, AMP desactivado y cero
imágenes faltantes. La consolidación utiliza exclusivamente los archivos
guardados; no ejecuta otra inferencia ni vuelve a evaluar test.

\subsection{Evolución del entrenamiento}

F1 se expresa en la escala de 0 a 1. La exactitud se presenta como porcentaje
con símbolo explícito; la pérdida es entropía cruzada y el tiempo está en minutos.
\begin{center}\small
\begin{tabular}{rrrrrr}
\toprule
\textbf{Época} & \textbf{Pérdida train} & \textbf{F1 imagen val} & \textbf{F1 obs. val} & \textbf{Exactitud imagen} & \textbf{Minutos}\\
\midrule
1 & 1,0928 & 0,7883 & 0,8454 & 83,35\,\% & 97,2\\
2 & 0,8142 & 0,8430 & 0,9033 & 86,72\,\% & 47,4\\
3 & 0,7532 & 0,8620 & 0,9005 & 90,17\,\% & 47,9\\
4 & 0,7147 & 0,8436 & 0,8972 & 88,70\,\% & 48,1\\
5 & 0,6932 & 0,8680 & 0,9183 & 90,76\,\% & 48,2\\
6 & 0,6736 & 0,8564 & 0,9233 & 90,45\,\% & 54,9\\
7 & 0,6553 & 0,8725 & 0,9221 & 89,87\,\% & 52,6\\
8 & 0,6358 & 0,8782 & 0,9353 & 91,76\,\% & 60,3\\
9 & 0,6234 & 0,8825 & 0,9343 & 91,37\,\% & 52,7\\
10 & 0,6051 & 0,8803 & 0,9491 & 91,80\,\% & 47,7\\
11 & 0,5993 & 0,8997 & \textbf{0,9545} & 92,07\,\% & 94,6\\
12 & 0,5855 & 0,8924 & 0,9369 & 92,00\,\% & 47,4\\
13 & 0,5748 & 0,8788 & 0,9348 & 92,03\,\% & 47,3\\
14 & 0,5651 & 0,8992 & 0,9454 & 93,56\,\% & 47,2\\
15 & 0,5556 & 0,8948 & 0,9372 & 93,42\,\% & 47,3\\
16 & 0,5483 & 0,9121 & 0,9494 & 93,78\,\% & 47,3\\
\bottomrule\end{tabular}\end{center}

La pérdida disminuye de 1,0928 a 0,5483. La macro-F1 de observación pasa
de 0,8454 en la primera época a 0,9545 en la seleccionada (mejora de
10,92 puntos porcentuales en la escala de F1). La época 16 alcanza
0,9121 en F1 de imagen, pero 0,9494 en F1 de observación; por ello no
sustituye al checkpoint de la época 11. La elección respeta el criterio
definido previamente y no utiliza el rendimiento de test.

La suma de los tiempos registrados para las 16 épocas es de 14,80 horas;
la media es 55,5 minutos por época.
Estos tiempos incluyen validación dentro de cada época, pueden reflejar
suspensiones y no equivalen al tiempo continuo entre el inicio y el fin
de una ejecución con pausas y reanudaciones.


\begin{figure}[htbp]
\centering
\begin{tikzpicture}
\begin{axis}[width=.94\textwidth,height=5.5cm,xlabel={Época},ylabel={Entropía cruzada},
xmin=1,xmax=16,xtick={1,4,8,11,16},grid=major,
tick label style={font=\small},label style={font=\small}]
\addplot[blue,mark=*,mark size=1.5pt] coordinates {(1,1.09280272) (2,0.81420936) (3,0.75321910) (4,0.71470973) (5,0.69315276) (6,0.67358492) (7,0.65531671) (8,0.63580538) (9,0.62341348) (10,0.60508928) (11,0.59928756) (12,0.58554775) (13,0.57480244) (14,0.56513808) (15,0.55558786) (16,0.54828336)};
\addplot[red,dashed] coordinates {(11,0.54)(11,1.1)};
\end{axis}
\end{tikzpicture}
\par\smallskip
\begin{tikzpicture}
\begin{axis}[width=.94\textwidth,height=6cm,xlabel={Época},ylabel={Macro-F1 (0--1)},
xmin=1,xmax=16,ymin=.77,ymax=.97,xtick={1,4,8,11,16},grid=major,
legend pos=south east,legend style={font=\small},tick label style={font=\small},label style={font=\small}]
\addplot[teal,mark=*,mark size=1.5pt] coordinates {(1,0.84537213) (2,0.90328831) (3,0.90047605) (4,0.89724439) (5,0.91827070) (6,0.92326025) (7,0.92208221) (8,0.93527384) (9,0.93433173) (10,0.94907896) (11,0.95454277) (12,0.93685406) (13,0.93480094) (14,0.94544569) (15,0.93718054) (16,0.94938297)};
\addlegendentry{Val por observación}
\addplot[blue,mark=square*,mark size=1.5pt] coordinates {(1,0.78831466) (2,0.84302836) (3,0.86196481) (4,0.84363064) (5,0.86797921) (6,0.85643758) (7,0.87253452) (8,0.87815451) (9,0.88246991) (10,0.88031467) (11,0.89968718) (12,0.89241032) (13,0.87879134) (14,0.89923725) (15,0.89483876) (16,0.91210182)};
\addlegendentry{Val por imagen}
\addplot[red,dashed,forget plot] coordinates {(11,.77)(11,.97)};
\end{axis}
\end{tikzpicture}
\caption{Curvas construidas a partir del historial guardado. La línea roja señala
la época 11; las cinco épocas siguientes no mejoran el criterio de selección.
El eje de F1 está restringido a 0,77--0,97 para mostrar su evolución.}
\end{figure}

\subsection{Evaluación final del checkpoint seleccionado}

\begin{center}\small\begin{tabular}{llrrrr}
\toprule
\textbf{Conjunto} & \textbf{Unidad} & \textbf{N} & \textbf{Macro-F1} & \textbf{Exactitud} & \textbf{Top-3}\\
\midrule
VAL & Imagen & 6\,290 & 0,8997 & 92,07\,\% & 98,30\,\%\\
VAL & Observación & 3\,452 & 0,9545 & 95,94\,\% & 99,33\,\%\\
TEST & Imagen & 7\,159 & 0,8875 & 91,09\,\% & 98,42\,\%\\
TEST & Observación & 4\,039 & 0,9463 & 95,79\,\% & 99,43\,\%\\
\bottomrule\end{tabular}\end{center}

La macro-F1 de observación es 0,9545 en val y 0,9463 en test; la
diferencia es de 0,82 puntos porcentuales en la escala de F1. La exactitud
por observación en test es 95,79\,\%. F1 combina precisión y sensibilidad
y no representa directamente el porcentaje de aciertos.

En val se registran 499 errores de imagen y 140 de observación; en test,
638 y 170, respectivamente. El promedio de probabilidades de las fotografías
de un hallazgo permite una predicción conjunta con mayor rendimiento
observado que la evaluación de fotografías individuales. Las unidades y
sus denominadores son diferentes; estas diferencias no constituyen una
estimación causal del aporte de cada fotografía.

\subsection{Rendimiento por género y soporte}

\begin{center}\small\begin{tabular}{lrrrrrr}\toprule
& \multicolumn{3}{c}{\textbf{Val: observación}} & \multicolumn{3}{c}{\textbf{Test: observación}}\\
\textbf{Género} & \textbf{N} & \textbf{F1} & \textbf{Recall} & \textbf{N} & \textbf{F1} & \textbf{Recall}\\
\midrule
Amanita & 904 & 0,9804 & 96,90\,\% & 1\,315 & 0,9797 & 97,34\,\%\\
Auricularia & 288 & 0,9589 & 93,06\,\% & 319 & 0,9477 & 90,91\,\%\\
Ganoderma & 335 & 0,9273 & 91,34\,\% & 346 & 0,9283 & 91,62\,\%\\
Geastrum & 79 & 0,9325 & 96,20\,\% & 182 & 0,9508 & 95,60\,\%\\
Hericium & 66 & 0,9219 & 89,39\,\% & 51 & 0,9184 & 88,24\,\%\\
Lycoperdon & 408 & 0,9377 & 97,79\,\% & 407 & 0,9382 & 98,77\,\%\\
Morchella & 11 & 1,0000 & 100,00\,\% & 20 & 0,9474 & 90,00\,\%\\
Pleurotus & 251 & 0,9569 & 97,21\,\% & 319 & 0,9484 & 94,98\,\%\\
Trametes & 756 & 0,9649 & 96,30\,\% & 791 & 0,9639 & 96,33\,\%\\
Xylaria & 354 & 0,9650 & 97,46\,\% & 289 & 0,9408 & 96,19\,\%\\
\bottomrule\end{tabular}\end{center}

El archivo de métricas por género conserva también precisión, F1 y
sensibilidad a nivel de imagen, junto con verdaderos positivos, falsos
positivos y falsos negativos para los dos conjuntos.

\subsection{Matrices de confusión}

Las filas representan el género real y las columnas el predicho; todos
los valores de las siguientes matrices son \textbf{conteos enteros de observaciones}.
Las columnas se abrevian en el mismo orden de las clases: Aman, Auri,
Gano, Geas, Heri, Lyco, Morc, Pleu, Tram y Xyla.
\begin{center}\scriptsize
\setlength{\tabcolsep}{3pt}
\begin{tabular}{lrrrrrrrrrr}\toprule
\multicolumn{11}{c}{\textbf{VAL: matriz por observación}}\\
Real / Predicho & Aman & Auri & Gano & Geas & Heri & Lyco & Morc & Pleu & Tram & Xyla\\\midrule
Amanita & 876 & 0 & 0 & 3 & 1 & 22 & 0 & 0 & 1 & 1\\
Auricularia & 1 & 268 & 4 & 0 & 2 & 0 & 0 & 5 & 3 & 5\\
Ganoderma & 1 & 1 & 306 & 0 & 0 & 5 & 0 & 2 & 14 & 6\\
Geastrum & 0 & 0 & 0 & 76 & 0 & 2 & 0 & 1 & 0 & 0\\
Hericium & 0 & 0 & 0 & 0 & 59 & 1 & 0 & 1 & 5 & 0\\
Lycoperdon & 4 & 0 & 2 & 3 & 0 & 399 & 0 & 0 & 0 & 0\\
Morchella & 0 & 0 & 0 & 0 & 0 & 0 & 11 & 0 & 0 & 0\\
Pleurotus & 1 & 1 & 1 & 2 & 0 & 0 & 0 & 244 & 2 & 0\\
Trametes & 0 & 1 & 10 & 0 & 0 & 9 & 0 & 4 & 728 & 4\\
Xylaria & 0 & 0 & 2 & 0 & 0 & 5 & 0 & 2 & 0 & 345\\
\bottomrule\end{tabular}\end{center}
\begin{center}\scriptsize
\setlength{\tabcolsep}{3pt}
\begin{tabular}{lrrrrrrrrrr}\toprule
\multicolumn{11}{c}{\textbf{TEST: matriz por observación}}\\
Real / Predicho & Aman & Auri & Gano & Geas & Heri & Lyco & Morc & Pleu & Tram & Xyla\\\midrule
Amanita & 1280 & 0 & 0 & 6 & 0 & 26 & 0 & 2 & 1 & 0\\
Auricularia & 2 & 290 & 4 & 1 & 1 & 1 & 0 & 6 & 4 & 10\\
Ganoderma & 5 & 2 & 317 & 0 & 0 & 4 & 0 & 2 & 13 & 3\\
Geastrum & 1 & 0 & 0 & 174 & 0 & 5 & 0 & 1 & 0 & 1\\
Hericium & 0 & 0 & 0 & 0 & 45 & 1 & 0 & 1 & 4 & 0\\
Lycoperdon & 3 & 0 & 0 & 0 & 1 & 402 & 0 & 1 & 0 & 0\\
Morchella & 1 & 0 & 0 & 0 & 0 & 0 & 18 & 0 & 1 & 0\\
Pleurotus & 5 & 1 & 1 & 1 & 0 & 2 & 0 & 303 & 3 & 3\\
Trametes & 1 & 0 & 11 & 0 & 0 & 7 & 0 & 3 & 762 & 7\\
Xylaria & 0 & 0 & 4 & 2 & 0 & 2 & 0 & 1 & 2 & 278\\
\bottomrule\end{tabular}\end{center}

Se exportaron además matrices gráficas de imagen y observación, tanto
con conteos como normalizadas por fila. En estas últimas cada celda
indica explícitamente un porcentaje del soporte del género real.

\subsection{Análisis de errores de validación}

Los errores de val se reconstruyeron a partir de las predicciones guardadas
sin utilizar el modelo. Las probabilidades exportadas tienen seis decimales;
la agregación reprodujo exactamente la matriz de observación guardada.

\begin{center}\small\begin{tabular}{llrr}\toprule
\textbf{Género real} & \textbf{Predicción} & \textbf{Errores obs.} & \textbf{Soporte real}\\\midrule
Amanita & Lycoperdon & 22 & 904\\
Ganoderma & Trametes & 14 & 335\\
Trametes & Ganoderma & 10 & 756\\
Trametes & Lycoperdon & 9 & 756\\
Hericium & Trametes & 5 & 66\\
\bottomrule\end{tabular}\end{center}

Hericium registra siete observaciones incorrectas de 66; cinco se
confunden con Trametes. Geastrum presenta tres errores de 79. Morchella
acierta las once observaciones de val, pero tiene cuatro imágenes
incorrectas de 21: el acierto por hallazgo no implica acierto en cada fotografía.

El recall de Morchella de 100\,\% tiene un intervalo Wilson aproximado
al 95\,\% de 74,12\,\%--100\,\%; Hericium, 79,69\,\%--94,77\,\%;
y Geastrum, 89,42\,\%--98,70\,\%. Estos intervalos describen la incertidumbre
del recall debida al soporte reducido, bajo un supuesto binomial; no son
intervalos de F1 ni una corrección de posibles dependencias entre observaciones.

La inspección visual de ejemplos de val muestra tomas parciales de
superficies, ejemplares deteriorados y una imagen de microscopía. Algunos
errores reciben probabilidades elevadas. Estos ejemplos orientan una
revisión taxonómica y de condiciones de captura, sin demostrar que la
etiqueta sea incorrecta ni modificarla automáticamente.

Se conservan listas completas de errores por imagen y observación, con
identificadores y probabilidades, y una galería local de ejemplos. El
análisis detallado para orientar los siguientes experimentos se limita
a val; las cifras de test se reportan como evaluación final descriptiva.

\subsection{Conservación de la línea base}

El checkpoint de la época 11, la configuración efectiva, el manifiesto
del dataset, el historial, las métricas y las predicciones finales se
copiaron a un archivo de conservación independiente de la carpeta de corrida.
La igualdad de los hashes SHA-256 entre origen y copia fue verificada,
y los archivos de la corrida original permanecen sin cambios.

La huella SHA-256 del checkpoint conservado es:
\begin{center}\small\ttfamily
a5edc643662266708263899c2fdd6977f0dabac160739ba7ccfde82603af50e3
\end{center}

La configuración efectiva es la registrada durante el entrenamiento.
La conservación incluye hashes de cada archivo y permite utilizar esta
línea base como referencia para S5 y S6.
\section{Limitaciones del experimento}

\begin{itemize}
\item \textbf{Representatividad:} FungiTastic procede principalmente del contexto
danés. El rendimiento en sus particiones no valida automáticamente el desempeño
en fotografías colombianas o capturadas mediante móviles.
\item \textbf{Desbalance y soporte:} conservar el protocolo mantiene clases poco
representadas; macro-F1 y métricas por clase deben leerse junto con sus conteos.
\item \textbf{Dependencia visual:} una pareja perceptualmente similar permanece
entre train y test. Cero duplicados exactos o cruces sitio-día no descarta otros
tipos de correlación o errores de etiqueta.
\item \textbf{Resolución:} las imágenes descargadas son 300p y la red observa
recortes de 224 píxeles; puede perderse detalle útil para especies próximas.
\item \textbf{Modalidades:} el clima está ausente y los otros metadatos no se
usan en este modelo. Los resultados no corresponden a un modelo multimodal.
\item \textbf{Licencias:} el manifiesto documenta declaraciones inconsistentes
en FungiTastic y registra conservadoramente CC-BY-NC-SA para el dataset.
\end{itemize}

\section{Trabajo próximo a realizar: semanas S5 y S6}

La arquitectura de referencia finalizó su entrenamiento y evaluación, y el
checkpoint seleccionado quedó conservado junto con la configuración efectiva.
Esta línea base constituye la referencia para el siguiente trabajo de la Fase 3.

\subsection{S5: 5--11 de octubre. Arquitectura eficiente para móvil}

La semana S5 corresponde al entrenamiento de una arquitectura eficiente apta
para móvil, al ajuste de hiperparámetros y al manejo del desbalance.
Se plantea:
\begin{itemize}
\item Seleccionar y entrenar una arquitectura eficiente, conservando el dataset
v002 y el protocolo de evaluación para facilitar la comparación con ResNet-18.
\item Ajustar tasa de aprendizaje, regularización y otros hiperparámetros
mediante resultados de val, manteniendo test fuera de la selección.
\item Analizar estrategias de manejo del desbalance en train, con atención a
los géneros de menor soporte y a sus efectos sobre macro-F1 y sensibilidad.
\item Caracterizar tamaño del modelo, latencia de inferencia y viabilidad de
ejecución móvil; distinguir estas medidas del rendimiento de clasificación.
\end{itemize}

\subsection{S6: 12--18 de octubre. Consolidación y modularización}

La semana S6 corresponde a la consolidación de métricas base por nivel
taxonómico y a la modularización del pipeline para la Fase 5.
Se plantea:
\begin{itemize}
\item Consolidar métricas, matrices de confusión y análisis de errores de las
arquitecturas de referencia y eficiente, indicando soporte por clase y unidad
de evaluación (imagen u observación).
\item Organizar los resultados por nivel taxonómico. La versión v002 evalúa
géneros; cualquier evaluación a nivel de especie requerirá una definición
explícita de etiquetas, cobertura y particiones antes del entrenamiento.
\item Separar adquisición y curaduría, transformaciones, arquitectura,
entrenamiento y evaluación en módulos con interfaces y configuraciones comunes.
\item Preparar el pipeline para la integración prevista en la Fase 5,
conservando la procedencia de metadatos, el versionado y la trazabilidad experimental.
\end{itemize}

El resultado esperado de S5--S6 es una comparación reproducible entre desempeño
visual y requisitos de ejecución móvil, junto con una base modular para las
siguientes fases del proyecto. Las propuestas se evaluarán con val y las corridas
finales seleccionadas consultarán test una vez al concluir.

\section{Conclusión}

Se dispone de un dataset versionado y reproducible de diez géneros, con fotografías
descargadas, particiones oficiales y controles explícitos de integridad y fuga.
La primera línea base utiliza transferencia de aprendizaje con ResNet-18 y selección
por macro-F1 de observación en val. Se completaron 16 épocas y se seleccionó la
época 11; el test final alcanza macro-F1 de observación de 0,9463 y exactitud de
95,79\,\%. La revisión de errores de val y del soporte reducido de algunas clases
delimita la interpretación de estos resultados. La consolidación conservó el
checkpoint y su configuración sin repetir la evaluación de test. Las actividades de S3 y S4
establecen las condiciones metodológicas para el trabajo de S5 y S6: evaluar
una arquitectura eficiente para móvil, abordar el desbalance y consolidar un
pipeline modular con métricas claramente vinculadas al nivel taxonómico.

\end{document}
